% id: 0338
% book: RPM 중학 3-2 · 03 원과 직선
% section: 시험에 꼭 나오는 문제
% figure: crop:fig-0338.png
% answer: ③
% answer_source: 답지
% variant_level: 0 (원본 전사)
% 이 파일은 build.mjs 가 items.json 에서 만든다. 고칠 때는 items.json 을 고친다.
\smash{\raisebox{3.5mm}{\dmkichul{RPM 중학 3-2 0338번 · 59쪽}}}
\noindent\begin{minipage}[t]{0.58\linewidth}\vspace{0pt}\raggedright 오른쪽 그림과 같이 중심이 같은 두 원에서 큰 원의 현~$\pt{AB}$는 작은 원의 접선이고 점~$\pt{P}$는 그 접점이다. $\seg{AB}=6\sqrt{3}\mkern3mu \mathrm{cm}$, $\seg{OP}=3\mkern3mu \mathrm{cm}$일 때, 큰 원의 반지름의 길이는?\end{minipage}\hfill\begin{minipage}[t]{0.4\linewidth}\vspace{0pt}\centering \setlength{\dmfigw}{84.1pt}\ifdim\dmfigw>\linewidth\setlength{\dmfigw}{\linewidth}\fi\includegraphics[width=\dmfigw]{figures/fig-0338.png}\end{minipage}\par
\begin{choices32}{\rule[-3ex]{0pt}{8ex}$2\sqrt{7}\mkern3mu \mathrm{cm}$}{\rule[-3ex]{0pt}{8ex}$4\sqrt{2}\mkern3mu \mathrm{cm}$}{\rule[-3ex]{0pt}{8ex}$6\mkern3mu \mathrm{cm}$}{\rule[-3ex]{0pt}{8ex}$3\sqrt{5}\mkern3mu \mathrm{cm}$}{\rule[-3ex]{0pt}{8ex}$7\mkern3mu \mathrm{cm}$}\end{choices32}
